\chapter{Literature study and research gaps}
\label{ch:literature}

\section{Scope of the corpus}

The working library comprises five unique journal papers, the supporting information of \citet{Dobson2024}, and a duplicate file of \citet{Meenakshisundaram2020} that must not be cited twice.
Supporting classics (oxygen inhibition, reciprocity, Macosko--Miller gelation, CLIP, frontal photopolymerization) are used as they appear in those reports, not as an unbounded review of photopolymer science.

The five papers do not compete to be ``the PEGDA model.''
Each \emph{owns a layer} (Table~\ref{tab:papers} in Chapter~\ref{ch:background}).
A flowing grayscale hydrogel would need all layers, inherited from the paper that actually closed each one.

\section{High-fidelity kinetics: Dobson and Bowman}

\citet{Dobson2024} construct a first-principles description of multifunctional acrylate photopolymerization with a large species set: photoinitiator, primary radicals, short and network-bound macroradicals, unreacted and partially reacted monomer, oxygen, and terminated polymer.
Mobility is tied to free volume; termination includes reaction-diffusion control and chain-length effects; heat release is optional.
The illustrated chemistry is 1,6-hexanediol diacrylate (HDDA) with Irgacure~819, not PEGDA.

The campaign uses this paper as a \emph{reference ontology}, not as production rates.
The temptation to ``implement Dobson first in 2-D'' is the usual way a PhD project becomes un-debuggable: too many species, too many spatial operators, and no algebraic limit that can fail loudly.
If, after a correct Montgomery transcription, equal-dose FTIR still cannot be organised, Dobson's chain-length $k_t$ is the parked refinement---not the first commit.

\section{Production spatial kinetics: Montgomery et al.}

\citet{Montgomery2022} is the production constitutive law for Type~I PEGDA~$250$ / 0.1~wt\% Irgacure~819 in a thin film.
Four species are tracked: initiator $C_I$, radicals $C_R$, oxygen $C_O$, monomer (acrylate groups) $C_M$.
With diffusion,
\begin{equation}
  \partial_t C_\alpha = \nabla\cdot\bigl(D_\alpha(p)\nabla C_\alpha\bigr) + R_\alpha,
\end{equation}
and $R_\alpha$ given by their Eqs.~16--19.
Photolysis is first order in remaining initiator, with $\beta=2.70\times10^{-3}\,\mathrm{s^2\,kg^{-1}}$ multiplying $I$ in $\mathrm{W\,m^{-2}}$ so that $\beta I$ is a frequency.
Two radicals are produced per cleavage ($m=2$).
Termination is $2k_t C_R^2$ plus oxygen inhibition $k_O C_O C_R$.

The closures (their Eqs.~22--27, Table~2) make $k_p$ and $k_t$ conversion-dependent:
\begin{align}
  k_p(p) &= \frac{k_{p0}}{1+(k_{p0}/k_{p,D0})\exp(c_p p)}, \\
  k_t(p) &= \Bigl[\tfrac{1}{k_{t,\mathrm{SD}}}+\tfrac{\exp(c_p p)}{k_{t,\mathrm{TD0}}}\Bigr]^{-1}
           + C_{RD}(1-p)k_p(p).
\end{align}
(The paper writes $\exp(cp)$ for $\exp(c p)$ with $c$ the relative viscosity coefficient.)
Harmonic liquid--solid diffusivities (Eq.~28, Table~3) drop by four orders of magnitude from liquid to solid.
Beer--Lambert attenuation (Eqs.~4--6) uses molar initiator absorptivity, polymer versus monomer absorption, and an optional photoabsorber term $w A_{\mathrm{abs}}$.

The optical boundary condition (Eq.~7) superposes Gaussians of radius $\sigma=17.7\,\mu\mathrm{m}$ on a $50\,\mu\mathrm{m}$ pixel pitch, with $I_{\mathrm{pixel}}=41.08\,\mathrm{W\,m^{-2}}$ and $I_{\mathrm{ideal}}=64\,\mathrm{W\,m^{-2}}$.
G0 is full intensity and G100 is dark.
Their multilayer prints with 0.06~wt\% absorber are a different optical depth than the FTIR cell.

\paragraph{What we did not take from this paper as gospel.}
The Gaussian kernel is a fallback, not a measured DMD map.
The linear G0--G100 map used in code is \emph{assumed}, not their SI RGB calibration.
Table~2 rates are a fit to one thin-film recipe; they are not universal PEGDA constants.
Dissolved oxygen is not a Table~2 entry; the campaign flags $C_{O,0}$ as \stillreq.

\section{Projection kinematics: Meenakshisundaram, Sturm and Williams}

\citet{Meenakshisundaram2020} model scanning-mask projection vat photopolymerization.
The scientific object is the surface irradiance and the accumulated dose under a measured (generally non-Gaussian) pixel kernel, discrete frames, and a scan velocity.
They compare a Jacobs-style cure overlay to printed dimensions; that overlay is a diagnostic, not a species conversion.

The campaign Stage~3 implements conservation and bleed tests on Montgomery's Gaussian as a \emph{stand-in}.
A \tagmatch{} on Stage~3 requires a camera or radiometer map of the projector actually used, or digitization of their measured PSF.
Scan dose scaling as $1/v$ is a dwell identity for a translating compact pattern, not a statement that $p$ is independent of speed.

\section{Plug flow: Slutzky, Stone and Nunes}

\citet{Slutzky2019} pass a sheathed aqueous jet of PEGDA~$575$ (54~vol\%), water, dye, and 4~vol\% 2-hydroxy-2-methylpropiophenone through a $L=2.3\,\mathrm{mm}$ UV slit.
Steady plug flow gives
\begin{equation}
  U\frac{\mathrm{d}C_i}{\mathrm{d}x}=R_i\bigl(C,k_d(x)\bigr),\qquad
  k_d=0\text{ outside }[0,L].
\end{equation}
Their Eqs.~1a--d use $k_d=\phi\varepsilon I$ with $I$ in Einstein~$\mathrm{m^{-2}\,s^{-1}}$, termination $k_t R^2$ (no extra factor of two), and one radical per photolysis as written in their Table~1.
Taking the ratio of monomer and oxygen balances at constant $k_p,k_O$ yields Eq.~3,
\begin{equation}
  \frac{C_O}{C_{O,0}}=\Bigl(\frac{C_M}{C_{M,0}}\Bigr)^{k_O/k_p}.
\end{equation}
Gelation for comparison to fiber onset is the operational predicate $M/M_0=0.98$ (about $2\%$ double bonds reacted), following microfluidic convention \citep{Dendukuri2008}---not Zhu's $P_{\mathrm{gel}}$.
Oxygen diffusion into the $40\,\mu\mathrm{m}$ jet during residence is argued small relative to advection of the dissolved inventory.

Critical speed scales as $U_c\sim L k_d [\mathrm{PI}]_0/[\mathrm{O}_2]_0$.
Pulsed light organises fiber length by $t_{\mathrm{uv}}/t_L$ with $t_L=L/U$.

This chemistry must not share a parameter book with Montgomery PEGDA~$250$/Irgacure.

\section{Network mechanics: Zhu et al.}

\citet{Zhu2020} couple Type~II eosin~Y / TEOA bleaching kinetics to acrylate conversion and then to a modified Macosko--Miller recursive model with intramolecular loops.
The campaign implements \emph{only} the network post-processor (their Eqs.~21, 26, 32--38).
Eosin bleaching ODEs (their Eqs.~3--8) remain off unless a Type~II kinetics branch is explicitly opened.

The intermolecular probability is
\begin{equation}
  \theta=\frac{[\mathrm{PEGDA}]_0}{[\mathrm{PEGDA}]_0+2C_0},
\end{equation}
with fitted $C_0=3.45$~vol\% (Table~A2).
Then $P_{\mathrm{gel}}=(1-\theta)/(2\theta)=C_0/[\mathrm{PEGDA}]_0$.
If $[\mathrm{PEGDA}]_0\le C_0$, the precursor never gels at $p=1$.
Without loops, $\eta=p^2$.
With loops, $\eta$ is assembled from probabilities that $m\ge 3$ carbon ends lead to the infinite network.
The phantom modulus is $G=\eta N_A[\mathrm{PEGDA}]_0 k_B T$, so $G/G_{\mathrm{ideal}}=\eta$.

\section{Supporting literature that the five papers rest on}

\paragraph{Oxygen.}
\citet{Decker1985} established the kinetic role of $\mathrm{O}_2$ as a radical scavenger.
\citet{Dendukuri2008} showed that PDMS walls replenish oxygen by diffusion and set a dead zone; that is a different experiment from Slutzky's jet.
\citet{Tumbleston2015} turned a persistent inhibited liquid layer into CLIP.

\paragraph{Reciprocity.}
\citet{Wydra2014} showed that BisGMA/TEGDMA does not obey $I\times t$ equivalence when bimolecular termination dominates.
\citet{Lin2019} organised related effects of oxygen, viscosity, and time-varying intensity.

\paragraph{Networks.}
\citet{Flory1941} and \citet{Macosko1976,Miller1976} are the gelation and post-gel recursive backbone.
\citet{Zhong2016} quantified how primary loops reduce elasticity.
\citet{Anseth1994} and \citet{Benson1959} frame diffusion-controlled termination.
\citet{Cabral2005} described light-driven polymerization fronts.
\citet{Goodner2002} coupled heat and mass transfer in thick films.
\citet{Dadashi2016} reviews Type~I versus Type~II initiation.
\citet{Cramer2001} marks thiol--ene as a different polymerization, parked here.

\section{Working curves, CLIP, and why they are not the polymer state}

Stereolithography practice has long used a working curve $C_d=D_p\ln(E/E_c)$ \citep{Jacobs1992book}.
$D_p$ is an effective optical penetration; $E_c$ is a critical exposure at which a gel-like depth is first observed for a given resin and process.
The construction is invaluable for machine calibration.
It collapses oxygen, termination order, grayscale bleed, and flow into two fitted numbers.
Those numbers change when the lamp spectrum, inhibitor, or speed changes, which is exactly the symptom that they are not constitutive.

Continuous liquid interface production \citep{Tumbleston2015} exploits a persistent oxygen-inhibited dead zone so that a part can be drawn continuously from a vat.
That dead zone is a \emph{diffusive} oxygen supply from a permeable window, closer in spirit to \citet{Dendukuri2008} than to \citet{Slutzky2019}, whose jet carries its oxygen inventory by advection.
A CLIP-style boundary condition is therefore a later Stage~4 variant, not a reason to replace $p$ with $C_d$.

Frontal photopolymerization \citep{Cabral2005} is another macroscopic signature of the same microscopic rates: an optically thick, exothermic film can propagate a conversion front.
Stage~4's Beer--Lambert front at $H=200\,\mu\mathrm{m}$ is a thin, isothermal cousin of that phenomenon, not a heat-release model \citep{Goodner2002}.

\section{Undermind synthesis versus what was coded}

A broader modelling note exists in the DFG Undermind workspace (unified photopolymerization roadmap and Stage~1 PEGDA implementation notes).
That synthesis is right about history-dependent kinetics, independent pattern and resin velocities, and the need for a loop-aware network closure.

It is not internally consistent as an \emph{execution order}.
One passage suggests implementing Dobson first, or making the first solver two-dimensional with a measured PSF.
The Stage~1 PEGDA notes, by contrast, specify a local four-species ODE, BDF integration, and Dobson as a check.

The code follows the Stage~1 notes and the local five-PDF ladder.
Undermind extras (Dendukuri walls, Wydra named tests, four in-house reactors, $\sim 80\,\mu\mathrm{m}$ flowing DMD patches) are treated as future \emph{tests and boundary conditions}, not as extra architecture in Stages~0--6.

\section{The gap closed by this campaign}

Before this repository, the group had papers, not a citable stack.
The specific gap was operational:

\begin{enumerate}
  \item No provenance-locked parameter books, so Montgomery $\beta$ could silently be used with Slutzky Einstein intensities.
  \item No local ODE with named gates (dark $p=0$, positivity, failed reciprocity, oxygen induction).
  \item No constant-rate fork to test Slutzky Eq.~3 as an identity rather than a fitted curve.
  \item No projection engine that is forbidden from reporting conversion.
  \item No 1-D reaction--diffusion with the \emph{same} cell RHS as the ODE.
  \item No plug-flow solver on Slutzky's stoichiometry.
  \item No network post-processor that cannot retune $k_p$.
\end{enumerate}

The remaining gap is \emph{quantitative match} under the same formulation (Chapter~\ref{ch:matching}), plus parked physics (heat, scattering, NS--gel, inverse DMD).
